% =====================================================================
% DROP-IN DRAFT — NOT YET IN paper.tex
% Two results: (A2) an exact closed form for zeta_1, and (A1) the
% asymptotic base of M*(n).
%
% VERIFICATION STATUS OF EVERY CLAIM BELOW
%   [V-CC]  verified by an independent path in the repo (Clifford
%           3-design + exact outcome enumeration), reproduces the
%           committed zeta_1 within MC spread at every n reached
%   [V-EX]  exact identity, verified numerically
%   [V-EN]  verified by direct 4^n enumeration over depolarized Haar
%           states, n <= 6, outside the repo
%   [PEND]  NOT yet checked against the committed-seed states — see the
%           PASS 5 prompt; do not submit while any [PEND] remains
%   [MEAS]  measured, not derived; stated as such in the text
%
% NOTATION: uses \operatorname{Tr}. Replace with the paper's macro.
% PLACEMENT: intended as a new subsection after 3.4. Inserting it as 3.5
%   renumbers the present 3.5 -> 3.6 and 3.6 -> 3.7. Cross-references to
%   the alpha transition and to Fig. 1 must be re-checked.
% =====================================================================


% ============================ BLOCK A =================================
% New subsection. Insert immediately after the present Sec. 3.4.
% ======================================================================

\subsection{The exact projection variance, and the base it fixes}
\label{sec:zeta1-closed-form}

The obstruction of Section~3.4 is not that \(\zeta_1\) has no closed form.
It is that the closed form is cubic in the Pauli expectations where the
ansatz was quadratic. We give it here, and it fixes the growth base that
Section~3.5 fits.

Write \(P_s\) for the Pauli string indexed by \(s \in \{I,X,Y,Z\}^n\), let
\(\operatorname{supp}(s)\) be the qubits on which \(s\) is not the
identity, and write \(\langle P_s\rangle = \operatorname{Tr}(\rho P_s)\).
Call an ordered pair \((s,s')\) \emph{compatible} when \(s\) and \(s'\)
agree on \(\operatorname{supp}(s)\cap\operatorname{supp}(s')\), and for a
compatible pair let \(s\triangle s'\) be the string equal to \(s\) on
\(\operatorname{supp}(s)\setminus\operatorname{supp}(s')\), equal to
\(s'\) on \(\operatorname{supp}(s')\setminus\operatorname{supp}(s)\), and
identity elsewhere.

\begin{proposition}
\label{prop:zeta1}
For the local Haar single-qubit shadow of Section~2.1 and any \(n\)-qubit
state \(\rho\),
\[
\mathbb{E}\!\left[\operatorname{Tr}(G\rho)^2\right]
= 4^{-n} \sum_{(s,s')\ \mathrm{compatible}}
3^{\,|\operatorname{supp}(s)\cap\operatorname{supp}(s')|}\,
\langle P_s\rangle \langle P_{s'}\rangle \langle P_{s\triangle s'}\rangle ,
\]
and \(\zeta_1 = \mathbb{E}[\operatorname{Tr}(G\rho)^2] -
\operatorname{Tr}(\rho^2)^2\).
\end{proposition}

\noindent
\emph{Proof sketch.} Write the single-qubit snapshot as
\(G^{(q)} = (I + 3\,\hat n_q\!\cdot\!\vec\sigma)/2\), so that
\(G = 2^{-n}\bigotimes_q (I + 3\,\hat n_q\!\cdot\!\vec\sigma)\) and
\(\operatorname{Tr}(G\rho) = 2^{-n}\sum_s 3^{|s|} c_s(\hat n)
\langle P_s\rangle\) with \(c_s = \prod_{q\in\operatorname{supp}(s)}
(\hat n_q)_{s_q}\). Parametrise \(\hat n_q = \varepsilon_q \hat m_q\)
with \(\hat m_q\) uniform and \(\varepsilon_q = \pm 1\) the outcome. The
outcome probability is itself linear in the \(\varepsilon_q\), so
squaring and summing over \(\varepsilon\) leaves exactly those triples
\((u,s,s')\) with \(\operatorname{supp}(u) = \operatorname{supp}(s)
\triangle \operatorname{supp}(s')\); every qubit then carries either
zero or two factors of \(\hat m\), never one or three. Averaging
\(\hat m_q\) over the sphere gives \(\delta/3\) on each qubit of
\(\operatorname{supp}(s)\cup\operatorname{supp}(s')\), which forces
agreement on the overlap and contributes
\(3^{-(|s|+|s'|-|s\wedge s'|)}\), cancelling against \(3^{|s|+|s'|}\)
and leaving \(3^{|s\wedge s'|}\). \(\square\)

The identity is exact for every state, entangled or not, and reduces to
the single-qubit law of Section~3.3: at \(n=1\) the four contributing
pairs give \(\tfrac14(1 + 3t^2 + t^2 + t^2) = \tfrac14 + \tfrac54 t^2\).
%% [V-EX] n=1 reduction. [V-CC] agrees to 1e-15 across the Bloch range.
An independent evaluation, by exact single-qubit Clifford 3-design
average and exact outcome enumeration rather than by the estimator of
Section~3.1, reproduces the sampled \(\zeta_1\) at every system size
within its Monte Carlo spread.
%% [V-CC] n=2..8, deviations 0.25%, 0.03%, -0.04%, -0.21%, all inside
%% the committed zeta1_rel_spread.

Two consequences follow, and they settle what Section~3.4 left open.

\paragraph{The rejected ansatz is exactly the diagonal.}
Restricting Proposition~\ref{prop:zeta1} to \(s = s'\) gives
\(4^{-n}\sum_s 3^{|s|}\langle P_s\rangle^2\), which is precisely the
weight-only quadratic form of Section~3.4 with \(c_w = 3^w/4^n\). The
residual that form leaves is therefore not an unidentified shortfall
but exactly the off-diagonal sum. Fitting free weight coefficients to
the diagonal recovers \(3^w/4^n\) to \(9\times10^{-16}\), and fitting
the same form to the exact \(\zeta_1\) across the nineteen-state family
of Section~3.4 leaves a maximum relative residual of \(0.122\),
reproducing the figure reported there.
%% [V-CC] c_w recovery 8.9e-16; ansatz residual 0.122 vs committed 0.123.

\paragraph{The diagonal sets the growth base.}
Two exact facts about the diagonal fix it. First,
\(\sum_s 3^{|s|} = (1 + 3\cdot 3)^n = 10^n\), since each qubit
contributes one identity of weight \(3^0\) and three Paulis of weight
\(3^1\).
%% [V-EX] verified at n=2,3,4: 100, 1000, 10000.
Second, for a Haar-random pure state depolarized at rate \(q\), the
mean squared Pauli expectation over non-identity strings is
\((1-q)^2/(2^n+1)\) exactly.
%% [V-EX][V-CC] measured/predicted ratio 1.0000 at every n reached.
Together these give
\[
\mathrm{DIAG}(n) \;\longrightarrow\; (1-q)^2 \left(\tfrac{5}{4}\right)^{\!n},
\]
because \(10^n/(4^n 2^n) = (5/4)^n\).
%% [V-CC] DIAG/(5/4)^n = 0.589, 0.794, 0.747, 0.796, 0.796, 0.809, 0.806
%%        (n=2..8) against (1-q)^2 = 0.81.
%% [V-EN] independent enumeration: 0.683, 0.729, 0.767, 0.784, 0.799 (n=2..6).

The off-diagonal sum is bounded. Its terms carry three non-identity
Pauli expectations, whose Haar average is \(2/((d+1)(d+2))\) with
\(d = 2^n\) because the two commuting factors multiply to the third; the
compatible pairs carry total weight \(16^n\), of which \(10^n\) is the
diagonal. Hence
\[
\mathrm{OFF}(n) \;\longrightarrow\; 2(1-q)^3 ,
\]
approached from below with corrections of order \(2^{-n}\) and
\((5/8)^n\).
%% [PEND] limit 2(1-q)^3 = 1.458 at q=0.1. NOT yet checked on the
%%        committed-seed states. [V-EN] outside the repo: measured OFF
%%        = 0.683, 0.924, 1.120, 1.247, 1.338 (n=2..6) against the
%%        finite-n form, ratio closing 1.442 -> 1.022.

Since \(\mathrm{OFF}\) is bounded while \(\mathrm{DIAG}\) grows
geometrically, the diagonal dominates, and
\[
\zeta_1(n) \;\longrightarrow\; (1-q)^2\left(\tfrac{5}{4}\right)^{\!n}
+ 2(1-q)^3 - (1-q)^4 .
\]
At \(q = 0.1\) this is \(0.81\,(1.25)^n + 0.802\), which matches the
sampled \(\zeta_1\) to \(1.2\%\) at \(n=8\) and \(0.5\%\) at \(n=9\).
%% [PEND] depends on the OFF limit above.

The second projection variance has base \(7\) per qubit, which is the
exact value at a maximally mixed marginal recorded in Section~2.5 and
which Haar-typical states approach; the sampled ratio
\(\zeta_2/7^n\) is flat at approximately \(1.0\) across the range.
%% [MEAS] the prefactor C_2 is measured, not derived.
Therefore
\[
M^*(n) = \frac{\zeta_2}{2\zeta_1}
\;\longrightarrow\; \frac{C_2}{2(1-q)^2}\left(\frac{28}{5}\right)^{\!n},
\qquad \frac{28}{5} = 5.6 .
\]

\paragraph{What the data says about the base.}
The window-dependence reported in Section~3.5 and Appendix~D is the
finite-size approach to this constant, not an artifact of the fit. Three
diagnostics agree, none of which depends on a choice of fitting window:

\begin{center}
\begin{tabular}{lccc}
\hline
quantity & predicted & measured & window \\
\hline
\(\zeta_1\) geometric-mean ratio & \(5/4 = 1.25\)  & \(1.2433\) & \(n=4\ldots 9\) \\
\(\zeta_2\) geometric-mean ratio & \(7\)           & \(6.9732\) & \(n=4\ldots 9\) \\
\(M^*\)     geometric-mean ratio & \(28/5 = 5.6\)  & \(5.6087\) & \(n=4\ldots 9\) \\
\hline
\end{tabular}
\end{center}

\noindent
The \(M^*\) figure carries a \(95\%\) interval of \([5.457, 5.746]\),
which contains \(5.6\); \(M^*(n)/5.6^n\) is flat over \(n \ge 4\),
running \(0.557\) at \(n=4\) to \(0.562\) at \(n=9\). The base is
therefore derived rather than fitted. The prefactor is not: it depends
on \(C_2\), which we measure and do not derive, and the finite-size
values sit below the asymptotic \(C_2/2(1-q)^2 \approx 0.62\) because
the constant term in \(\zeta_1\) is still \(12\%\) of it at \(n=9\).
%% Keep this caveat. It is the sentence a referee will test.


% ============================ BLOCK B =================================
% Patch to the closing paragraph of Sec. 3.4. The present text hedges
% the missing ingredient as "an observation from the algebra rather than
% a theorem" and declines to rule out a closed form. Both are now
% superseded. REPLACE that paragraph with:
% ======================================================================

A weight-only quadratic form is insufficient, and
Section~\ref{sec:zeta1-closed-form} identifies exactly what it omits: it
is the diagonal of an exact cubic identity, and the residual it leaves
is the off-diagonal sum of that identity. The shortfall is therefore not
an unresolved feature of the algebra. For \(n \ge 2\) the projection
variances are still estimated numerically in what follows, since the
exact form requires the full Pauli spectrum; the variance law is exact
and its inputs, as evaluated here, are sampled.


% ============================ BLOCK C =================================
% Patch to the conclusion. The present text opens the three questions
% with "Whether zeta_1 admits a closed form under a richer ansatz than
% Pauli weight, which would make M*(n) fully analytic." REPLACE that
% clause with:
% ======================================================================

Two questions follow from this work, the first of the three we opened
having been closed in Section~\ref{sec:zeta1-closed-form}: \(\zeta_1\)
does admit a closed form, cubic rather than quadratic in the Pauli
expectations, and it fixes the growth base of \(M^*(n)\) at \(28/5\).
What it does not fix is the prefactor, which carries the second
projection variance's constant and which we measure rather than derive.


% ============================ BLOCK D =================================
% MENTOR'S CALL — do not apply without him. These change the paper's
% headline framing, which the project has held fixed.
% ======================================================================

% D1. ABSTRACT. Present: "...this budget grows approximately
% exponentially, with a fitted base near 5.3 for this ensemble — an
% empirical finite-size fit over that range, not a proven asymptotic
% rate."
%
% Candidate replacement:
%   "...this budget grows exponentially with base 28/5, which we derive
%   from an exact closed form for the first projection variance and
%   which the finite-size fits approach from below over the range we
%   reach."
%
% Weigh against: the whole paper is framed as finite-size theory plus
% simulation. Promoting a derived constant into the abstract changes
% what the paper claims to be. Consider instead leaving the abstract
% alone and letting Sec. 3.5 carry it.

% D2. SEC 3.5 / APPENDIX D. The sentence "This is a finite-size fit over
% the sizes we reach ... not a measurement of an asymptotic rate" is
% still true of the FIT and should stay. What can be added is a forward
% reference: the window drift is the approach to the derived base of
% Sec. \ref{sec:zeta1-closed-form}, not an instability in the fit.
% This is the minimal honest change and costs nothing framing-wise.
